\documentclass[10pt,a4paper]{amsart}
\usepackage[margin=19mm]{geometry}
\usepackage{amsmath,amssymb,mathtools}
\usepackage[scr=rsfs,cal=euler]{mathalfa}
\usepackage{xfrac}
\usepackage[unicode,hidelinks,bookmarks=false]{hyperref}
\newcommand{\ie}{i.e.\ }
\newcommand{\cf}{cf.\ }
\newcommand{\resp}{respectively\ }
\newcommand{\Subsection}{\subsection}
\providecommand{\nobreakdash}{\nobreak\hbox{-}}
\setlength{\emergencystretch}{2em}
\multlinegap0pt
\usepackage{bm}
\usepackage{bbm}
\usepackage{dsfont}
\usepackage{xspace}
\usepackage{cancel}
\usepackage{mathtools}
\usepackage{amsthm}
\makeatletter
\@ifundefined{theorem}{\newtheorem{theorem}{Theorem}}{}
\@ifundefined{assumption}{\newtheorem{assumption}[theorem]{Assumption}}{}
\@ifundefined{definition}{\newtheorem{definition}[theorem]{Definition}}{}
\@ifundefined{proposition}{\newtheorem{proposition}[theorem]{Proposition}}{}
\makeatother
\DeclareMathOperator{\Fr}{Fr}
\newcommand{\Atom}{\mathrm{Atom}}
\newcommand{\Cn}{\mathrm{Cn}}
\title[Closure-Preserving Rate-Distortion for Reversible Logging]{Closure-Preserving Rate-Distortion for Reversible Logging}
\author{Jianfeng Xu}
\date{Source: arXiv:2606.16592v2 (CC BY 4.0)}
\begin{document}
\maketitle

\noindent\emph{Source and licence.} Closure-Preserving Rate-Distortion for Reversible Logging. Version arXiv:2606.16592v2 (CC BY 4.0), distributed under the Creative Commons Attribution 4.0 International licence, as recorded in the arXiv licence metadata for this exact version and shown on the abstract page https://arxiv.org/abs/2606.16592v2. Full rights notice: licenses/arxiv-2606.16592-CC-BY-4.0.txt.\par\smallskip

\noindent\textit{Editorial scope.} Counted unit: the Proposition whose source label is \texttt{prop:frontier-core-downward} in the article (source0.tex lines 1769--1783, source label \texttt{prop:frontier-core-downward}), delivered together with the article's own complete original proof of it (source0.tex lines 1785--1848, one \texttt{proof} environment of 64 lines). No mathematics is written, repaired or re-derived: statement and proof are the article's own text line for line. Required context retained uncounted: def:core (lines 660--666); assm:rcn-acyclic (lines 1763--1767). Every definition, display and label of those lines is retained unchanged. Omitted from this package and not redistributed: 1--659, 667--1762, 1768--1768, 1784--1784, 1849--2842. Nothing inside a retained excerpt is omitted or altered: every retained body is delivered exactly as the article's lines are written, so no omission receipt of any kind accompanies this package. Citations: the retained excerpts carry no citation command at all, so no bibliography entry of the article is transcribed and no reference is redistributed. Reference adaptation: none was needed and none was made; every reference inside the delivered unit resolves inside it, so no unresolved reference or citation is produced. Printed slips kept literal: no printed word, symbol, hypothesis or number of the counted statement or of its proof was altered. Layout and class: the article's own class is \texttt{article} with packages including amsmath, amssymb, amsthm, mathtools, graphicx, hyperref, geometry, array, fontenc, inputenc, tikz, pgfplots, booktabs, enumitem; the wrapper is the lane's pinned \texttt{amsart} editorial wrapper at 10pt on A4, which restates the theorem-like environments the retained text needs and adds the article's own macro definitions that the retained text needs (\texttt{\textbackslash Atom}, \texttt{\textbackslash Cn}, \texttt{\textbackslash Fr}), with extra wrapper packages (bm, bbm, dsfont, xspace, cancel, mathtools). The delivered numbering is the unit's own. Assets: no retained span contains a figure environment, a graphics-inclusion command, a PDF-page inclusion, a table or a TikZ picture; no image file is copied into this package, the package declares no asset dependency and no image file is redistributed with it. Licence: version arXiv:2606.16592v2 of this article is distributed under the Creative Commons Attribution 4.0 International licence, as recorded in the arXiv licence metadata for this exact version and shown on the abstract page https://arxiv.org/abs/2606.16592v2. A full rights notice is delivered at licenses/arxiv-2606.16592-CC-BY-4.0.txt. Characters: every retained excerpt is pure ASCII, so the delivered engine, XeLaTeX with its default Latin Modern text font, reports no missing character.
\begin{definition}[Irredundant core via deletion scan]
\label{def:core}
Define \(A=\Atom_{\mathrm{rev}}(S_O)\) by the deterministic deletion scan:
start from \(A\leftarrow S_O\); scan \(s\in S_O\) in \(\preceq\)-order and delete \(s\) if
\(s\in \Cn_{\mathrm{rev}}(A\setminus\{s\})\).
Let \(J\coloneqq S_O\setminus A\).
\end{definition}

\begin{assumption}[Acyclicity of causal reachability on configurations]
\label{assm:rcn-acyclic}
For every reachable configuration \(X\subseteq\overline{T}\), the restriction of
\(\mathsf{Cause}^\star\) to \(X\) is acyclic.
\end{assumption}

\begin{proposition}[Frontier core for disciplines with downward-closed presence]
\label{prop:frontier-core-downward}
Fix a configuration \(X\subseteq\overline{T}\), let
\[
S_O=S_O^{\mathrm{rcn}}(X), \qquad F\coloneqq \Fr(X), \qquad S_F\coloneqq S_O^{\mathrm{rcn}}(F).
\]
Assume Assumption~\ref{assm:rcn-acyclic}, that rule \emph{(CC1)} belongs to the closure system,
and that no rule other than \emph{(CC1)} derives atoms of the form \(\mathsf{In}(\cdot)\).
Then
\[
\Atom_{\mathrm{rev}}(S_O)=S_F.
\]
In particular, the identity holds for both \(\Cn^{\mathrm{causal}}_{\mathrm{rcn}}\) and
\(\Cn^{\mathrm{cr}}_{\mathrm{rcn}}\).
\end{proposition}

\begin{proof}
We first show that
\[
\Cn_{\mathrm{rev}}(S_F)=\Cn_{\mathrm{rev}}(S_O).
\]
Since \(S_F\subseteq S_O\), monotonicity gives
\[
\Cn_{\mathrm{rev}}(S_F)\subseteq \Cn_{\mathrm{rev}}(S_O).
\]
For the reverse inclusion, let \(t\in X\).

If \(t\in F\), then \(\mathsf{In}(t)\in S_F\subseteq \Cn_{\mathrm{rev}}(S_F)\).

If \(t\notin F\), then \(t\) has a strict \(\mathsf{Cause}^\star\)-successor in \(X\).
Because \(X\) is finite and \(\mathsf{Cause}^\star\!\upharpoonright_X\) is acyclic, repeated choice of a strict
\(\mathsf{Cause}^\star\)-successor must terminate at a \(\mathsf{Cause}^\star\)-maximal element \(f\in F\) such that
\(\mathsf{Cause}^\star(t,f)\).
Since \(\mathsf{In}(f)\in S_F\), rule \emph{(CC1)} yields
\[
\mathsf{In}(f)\wedge \mathsf{Cause}^\star(t,f)\Rightarrow \mathsf{In}(t),
\]
so \(\mathsf{In}(t)\in \Cn_{\mathrm{rev}}(S_F)\).

Hence \(S_O\subseteq \Cn_{\mathrm{rev}}(S_F)\), and monotonicity plus idempotence imply
\[
\Cn_{\mathrm{rev}}(S_O)\subseteq \Cn_{\mathrm{rev}}(\Cn_{\mathrm{rev}}(S_F))=\Cn_{\mathrm{rev}}(S_F).
\]
Therefore \(\Cn_{\mathrm{rev}}(S_F)=\Cn_{\mathrm{rev}}(S_O)\).

Next we show that each frontier atom is irredundant.
Fix \(f\in F\) and suppose, towards a contradiction, that
\[
\mathsf{In}(f)\in \Cn_{\mathrm{rev}}\bigl(S_F\setminus\{\mathsf{In}(f)\}\bigr).
\]
By assumption, the only rule that derives new \(\mathsf{In}(\cdot)\) atoms is \emph{(CC1)}.
Consequently, any derivation of \(\mathsf{In}(f)\) from \(S_F\setminus\{\mathsf{In}(f)\}\) must
involve an instance of (CC1) with conclusion \(\mathsf{In}(f)\).  The premises of (CC1) require
some \(\mathsf{In}(g)\) together with \(\mathsf{Cause}^\star(f,g)\).  Since the derivation
uses only facts from \(S_F\setminus\{\mathsf{In}(f)\}\) and Horn rules, the atom
\(\mathsf{In}(g)\) must eventually be supported by a ground fact in
\(S_F\setminus\{\mathsf{In}(f)\}\); that is, there exists
\(\mathsf{In}(g)\in S_F\setminus\{\mathsf{In}(f)\}\) such that
\(\mathsf{Cause}^\star(f,g)\).  But then \(g\in X\) is a strict
\(\mathsf{Cause}^\star\)-successor of \(f\), contradicting the maximality of \(f\in F\).
Hence
\[
\mathsf{In}(f)\notin \Cn_{\mathrm{rev}}\bigl(S_F\setminus\{\mathsf{In}(f)\}\bigr)
\qquad\text{for all } f\in F.
\]

Finally, consider the deletion scan of Definition~\ref{def:core}.
By the previous paragraph, no frontier atom is ever deleted.
When a non-frontier atom \(\mathsf{In}(t)\) is scanned, the current working set still contains \(S_F\), and from the first part
of the proof we already know
\[
\mathsf{In}(t)\in \Cn_{\mathrm{rev}}(S_F).
\]
Thus
\[
\mathsf{In}(t)\in \Cn_{\mathrm{rev}}(A\setminus\{\mathsf{In}(t)\}),
\]
so \(\mathsf{In}(t)\) is deleted by the scan.
Therefore the scan keeps exactly the frontier atoms, and the final core is \(S_F\).
\end{proof}

\end{document}
